\documentclass[12pt]{exam}
\usepackage{amsthm}
\usepackage{libertine}
\usepackage[utf8]{inputenc}
\usepackage[margin=1in]{geometry}
\usepackage{amsmath,amssymb}
\usepackage{multicol}
\usepackage[shortlabels]{enumitem}
\usepackage{siunitx}
\usepackage{cancel}
\usepackage{graphicx}
\usepackage{pgfplots}
\usepackage{listings}
\usepackage{tikz}


\pgfplotsset{width=10cm,compat=1.9}
\usepgfplotslibrary{external}
\tikzexternalize

\newcommand{\class}{Math Research} % This is the name of the course 
\newcommand{\examnum}{HW3} % This is the name of the assignment
\newcommand{\examdate}{10/21/26} % This is the due date
\newcommand{\timelimit}{}





\begin{document}
\pagestyle{plain}
\thispagestyle{empty}

\noindent
\begin{tabular*}{\textwidth}{l @{\extracolsep{\fill}} r @{\extracolsep{6pt}} l}
\textbf{\class} & \textbf{Name:} & \textit{Toby Semple}\\ %Your name here instead, obviously 
\textbf{\examnum} &&\\
\textbf{\examdate} &&\\
\end{tabular*}\\
\rule[2ex]{\textwidth}{2pt}
% ---




\begin{itemize}
  \item[4.1:] Prove $2^{n}<n$ ! for $n \geq 4$.\\
  Base
$$
\begin{aligned}
& 2^{4}=16 \quad 4!=24 \\
& 16<24 \text { is true }
\end{aligned}
$$
Induction:
$$
=2^{2 n+1} \cdot 2=(n+1)!(n+1)
$$
\end{itemize}

$2^{n}$ only gets multiplied by 2 whereas n! get multiplied by n+1. Since the base case is 4, which is arendy bigger than 2), $(n+1)!$ is multiplied more than $2^{n+1}$ meaning $2^{n+1}$ will always be less than $(n+1)$ ! after the base case.

\begin{itemize}
  \item[4.3:] Prove $f_{0}+\ldots+f_{n}=f_{n+2}-1$ for the Fib. Seq.\\
Base case
$$
f_{0}=0, \quad f_{n+2}-1=f,-1=1-1=0
$$
\end{itemize}

Induction:

$$
\begin{aligned}
& \sum_{i=1}^{2} f_{i}=f_{0}+\ldots+f_{n}=f_{n+2}-1 \\
& \sum_{i=1}^{n+1} f_{i}=f_{n+3}-1=\left(f_{n}+f_{n}\right)=1 \\
& =\left(f_{n \times 2}-1\right)^{n+1} f_{n}=\frac{2}{i n}+f_{n}
\end{aligned}
$$

\begin{itemize}
  \item[4.5:] prove $\left(f_{0}\right)^{2}+\ldots+\left(f_{n}\right)^{2}=f_{n} f_{n+1}$ if seq start @ 0,2
$$
n=0: \quad 0^{2}=0.2
$$
  \item[] $n+1:\left(\sum_{i=0}^{n}\left(f_{n}\right)^{2}=f_{n} f_{n+1}\right)+\left(f_{n+1}\right)^{2}$
$$
\begin{aligned}
\sum_{i=0}^{n}\left(f_{n}\right)^{2}+\left(f_{n+1}\right)^{2}=\sum_{i=0}^{n+1}\left(f_{n+1}\right)^{2} & =f_{n+1}\left(f_{n}+f_{n+1}\right) \\
& =f_{n+1}^{n+1} f_{n+2}
\end{aligned}
$$
  \item[] 4.7:
$$
n=0: \quad 2^{n}=0 . \quad n=1: 2^{n}=2
$$
$n+1: \quad 2^{n+1}=2^{n} \cdot 2$
\end{itemize}

If set $A$ has $n$ elements and ( $2^{n}=a$ ) subsets, then set $B$ with $n+1$ elements would have $a \cdot 2$ subsets, - a copy of subsets ' $a$ ' and copy of $a$ ' containing $n + 1$.\\
4.9:

$$
n=1: \frac{1}{(12)}=1
$$

$$
n+1:\left(\sum_{i=0}^{n}\left(\overrightarrow{b_{1}}\right)=\frac{n}{(n+1)}\right)+\frac{1}{\left.m^{n}(n), 2 n\right)}
$$

$$
\begin{aligned}
\left(\sum_{i=0}^{n} \frac{1}{n(n+1)}\right)+\frac{1}{(n-1)(n+2)}=\sum_{i=0}^{n-10}\left(\frac{1}{n(n+1)}\right) & =\frac{n(n+2)+1}{(n+1)(n+2)} \\
& =\frac{(n+1)(n+2)}{n+2} \\
& =\frac{(n+1)}{n+1}
\end{aligned}
$$

\begin{itemize}
  \item[] 4.13:\\
True; there are no polka-dotted elephants, so there's no way to refute the claim.
  \item[] 4.15:\\
This 'proof' is false, as it trents infinity as a number instent of a symbol representing a concept.
  \item[] 4.17:\\
LNP is based upon the fact there's a 'lowest integer'.\\
\\
Strong Induction is also based upon the concept of a 'lowest integer'.\\
\\
Common Induction is based upon LNP, which means it's indirectly related to the concept of a 'lowest integer'.\\
\\
Since the truth whe of all of the primiples are baged upon the truth value of the same poinciple, the truth value of one priciple informs one of the truth value of the underlying concept, indirectly proving/ disproving the other two principles.
\end{itemize}


\end{enumerate}

Look! Math! $f(x)=4x^2 - \frac{3}{4}x +1$ %use the dollar signs to go into math mode


\end{document}